\chapter{MySQL Database Implementation}
All outputs in this chapter are real: they were produced on \tbl{hospital\_db\_test}, a database built from the seed
data and every migration (same schema as \tbl{hospital\_db}, checked by a schema-diff script). Statements that change
data were run inside a transaction that was then rolled back.

\section{Database Creation}
The schema is built by numbered migration files (\path{database/migrations/000} to \path{023}), applied in order by
\path{server/scripts/migrate.js}, which backs up the database first and records each file in
\tbl{schema\_migrations}. The database and the restricted account the application runs as:
\begin{lstlisting}[caption={Creating the database and the application's restricted account}]
CREATE DATABASE hospital_db CHARACTER SET utf8mb4 COLLATE utf8mb4_0900_ai_ci;
CREATE USER 'hospital_app'@'127.0.0.1' IDENTIFIED BY RANDOM PASSWORD;
GRANT SELECT, INSERT, UPDATE, DELETE, EXECUTE ON `hospital_db`.* TO 'hospital_app'@'127.0.0.1';
\end{lstlisting}
The application account cannot create, alter or drop tables or triggers; migrations use a separate administrative account.

\section{Table Creation}
Six key tables are shown below, exactly as MySQL reports them (\tbl{SHOW CREATE TABLE}); all 33 are in the complete
script (Appendix~A).

\lstinputlisting[caption={\tbl{patients}: unique ABHA number with a format CHECK}]{generated/create-patients.sql}
\lstinputlisting[caption={\tbl{encounters}: the visit, with its 1:1 link to an appointment and the parent key used by consistency keys}]{generated/create-encounters.sql}
\lstinputlisting[caption={\tbl{admissions}: five foreign keys with RESTRICT, SET NULL and a composite parent key}]{generated/create-admissions.sql}
\lstinputlisting[caption={\tbl{prescription\_items}: structured dosing and the dispensed CHECK}]{generated/create-prescription_items.sql}
\lstinputlisting[caption={\tbl{medicine\_batches}: a weak entity (UNIQUE owner + batch number)}]{generated/create-medicine_batches.sql}
\lstinputlisting[caption={\tbl{bill\_items}: the generated total and CASCADE from its bill}]{generated/create-bill_items.sql}

\section{Data Insertion}
\begin{lstlisting}[caption={Inserting a patient and a bill line (results below)}]
INSERT INTO patients (name, date_of_birth, gender, blood_group, phone, abha_number)
VALUES ('Meera Iyer', '1992-07-14', 'FEMALE', 'B+', '9810012345', '12345678901234');
INSERT INTO bill_items (bill_id, description, category, quantity, unit_price)
VALUES (1, 'Dressing change', 'PROCEDURE', 3, 150.00);
\end{lstlisting}
\textbf{Result:} 1 row each. The patient got \tbl{patient\_id = 37}. The bill line got \tbl{total\_price = 450.00}
without the statement supplying it (generated column). The seed data (\tbl{database/seed.sql}) inserts 17 staff
accounts, 12 doctors, 9 departments, 6 wards, 35 beds, 32 patients, 22 medicines, 15 lab tests and their activity.

\section{Data Modification and Deletion}
\begin{lstlisting}[caption={UPDATE and DELETE, including a cascade and a refused delete (results below)}]
UPDATE bill_items SET quantity = 4 WHERE item_id = 20;
UPDATE beds SET status = 'MAINTENANCE' WHERE bed_id = 1 AND status = 'AVAILABLE';
DELETE FROM bills WHERE bill_id = 1;          -- its bill_items go with it (ON DELETE CASCADE)
DELETE FROM patients WHERE patient_id = 1;    -- refused: the patient is referenced (RESTRICT)
\end{lstlisting}
\textbf{Results:} the bill line's total became 600.00 (4 $\times$ 150.00); bed GEN-A01 became MAINTENANCE (1 row);
deleting bill 1 removed 1 bill and its 8 items (8 $\to$ 0); deleting patient 1 failed with
\tbl{ER\_ROW\_IS\_REFERENCED\_2 (1451)}: foreign key \tbl{bed\_assignment\_log\_ibfk\_2} (\tbl{ON DELETE RESTRICT}).

\section{Sample Database Contents}
\begin{table}[H]\centering
\caption{Sample contents of \tbl{patients}, \tbl{departments} and beds per ward (\tbl{SELECT} on the test database)}\label{tab:sample}
{\scriptsize\setlength{\tabcolsep}{3pt}\begin{tabular}{lllll}\toprule
\textbf{patient\_id} & \textbf{name} & \textbf{date\_of\_birth} & \textbf{gender} & \textbf{blood\_group} \\\midrule
1 & Aarav Sharma & 1985-03-15 & MALE & O+ \\
2 & Diya Patel & 1990-07-22 & FEMALE & A+ \\
3 & Vihaan Reddy & 1978-11-30 & MALE & B+ \\
4 & Ananya Gupta & 1995-05-10 & FEMALE & AB+ \\
5 & Aditya Kumar & 1960-01-20 & MALE & O- \\
\bottomrule\end{tabular}}
\hfill{\scriptsize\setlength{\tabcolsep}{3pt}\begin{tabular}{lll}\toprule
\textbf{department\_id} & \textbf{name} & \textbf{floor} \\\midrule
1 & Emergency & 1 \\
2 & Cardiology & 2 \\
3 & Neurology & 3 \\
4 & Orthopedics & 2 \\
5 & General Medicine & 1 \\
6 & Pediatrics & 1 \\
\bottomrule\end{tabular}}
\hfill{\scriptsize\setlength{\tabcolsep}{3pt}\begin{tabular}{lll}\toprule
\textbf{ward} & \textbf{beds} & \textbf{available} \\\midrule
General Ward A & 10 & 8 \\
General Ward B & 8 & 7 \\
Cardiac Ward & 4 & 2 \\
Pediatric Ward & 3 & 2 \\
ICU Ward & 6 & 2 \\
Emergency Observation & 4 & 3 \\
\bottomrule\end{tabular}}

\end{table}

\section{SQL Query Implementation}
The fifteen queries below cover the ten required kinds, and most are the queries the application itself runs or
reports on. Each gives its SQL, purpose and real output.
\query{1}{Simple SELECT with conditions}{Medicines whose stock has fallen to or below their reorder level (the low-stock alert).}{q1}
\query{2}{ORDER BY}{Available ICU and ventilator beds, ordered by ward and bed number (what the emergency desk looks at first).}{q2}
\query{3}{Aggregates: COUNT, SUM, AVG, MIN, MAX}{Billing summary: how many bills, how much billed and paid, and the smallest, average and largest bill.}{q3}
\query{4}{GROUP BY with HAVING}{Departments that have more than one doctor, with their doctor count and combined capacity.}{q4}
\query{5}{INNER JOIN}{Today's appointments with patient and doctor names; the department comes from the doctor through the view v\_appointments.}{q5}
\query{6}{LEFT JOIN}{Every doctor with the number of ACTIVE admissions they hold, including doctors with none (LEFT JOIN keeps them, COUNT gives 0).}{q6}
\query{7}{Self join (another join type)}{Pairs of doctors who share a specialization: the doctors table joined to itself.}{q7}
\query{8}{Nested query (subquery in WHERE)}{Medicines priced above the average price of all medicines.}{q8}
\query{9}{Correlated subquery}{Each admitted patient with the time of their latest vital-signs reading (NULL: not taken yet). The subquery runs once per admission.}{q9}
\query{10}{IN and EXISTS}{Patients with an ACTIVE admission (EXISTS) who also have a lab order that is not yet completed (IN).}{q10}
\query{11}{ALL and ANY}{Medicines dearer than every antibiotic (> ALL), and antacids dearer than at least one antiplatelet (> ANY).}{q11}
\query{12}{Multiple tables (five-table join)}{Current inpatients with ward, bed, admitting doctor and department.}{q12}
\query{13}{Business question: bed occupancy}{Occupancy by bed type: how full is the hospital? (CASE inside SUM.)}{q13}
\query{14}{Business question: stock that expires soonest}{For each medicine with stock, its usable quantity and next expiry, from the batches (view v\_current\_stock).}{q14}
\query{15}{View}{The view current\_admissions\_view joins admissions with patients, doctors, beds, wards and departments; callers select from it like a table.}{q15}


\section{Advanced DBMS Features}
\textbf{Views (10).} \tbl{v\_appointments} returns appointments with the department taken from the doctor
(Section~6.3), so reads keep their columns after the decomposition; \tbl{v\_current\_stock} shows each medicine's
total, usable and expired quantities and next expiry from its batches (Query~14); \tbl{current\_admissions\_view}
(Query~15), \tbl{available\_beds\_view}, \tbl{bed\_status\_summary\_view}, \tbl{hospital\_occupancy\_view},
\tbl{available\_doctors\_view}, \tbl{doctor\_workload\_view}, \tbl{emergency\_queue\_view} and
\tbl{pending\_bills\_view} give reporting summaries.

\textbf{Triggers (7).} \tbl{after\_admission\_insert} (Listing~\ref{lst:trigger}) and \tbl{after\_admission\_discharge}
keep beds, doctor workload and the bed history in step with admissions;
\tbl{after\_prescription\_dispense} takes a dispensed quantity off \tbl{medicines.stock\_quantity};
\tbl{before\_medicine\_stock\_update} refuses a negative total; \tbl{audit\_logs\_no\_update} and
\tbl{audit\_logs\_no\_delete} make the audit log append-only; \tbl{after\_bed\_assignment\_log} is an empty
placeholder from the original schema.
\lstinputlisting[caption={Trigger \tbl{after\_admission\_insert} (since migration 020, workload counts only ACTIVE admissions)},label={lst:trigger}]{generated/trigger-admission.sql}

\textbf{Stored procedures (5)}, from the original schema: \tbl{sp\_admit\_patient}, \tbl{sp\_discharge\_patient},
\tbl{sp\_allocate\_emergency}, \tbl{sp\_generate\_bill} and \tbl{sp\_get\_emergency\_queue}.

\textbf{Transactions and locking.} Every change runs in a transaction (\tbl{withTransaction} in
\tbl{server/utils/audit.js}): the change and its audit row commit together or not at all. InnoDB runs at REPEATABLE
READ, and operations that read a value and then change it lock the rows first (\tbl{SELECT \dots{} FOR UPDATE}).
To prevent deadlocks, rows are always locked in the same order. Dispensing, for example, locks the prescription,
then its medicines in ascending id, then each medicine's batches by expiry date:
\begin{lstlisting}[caption={Lock order in one dispensing transaction (\tbl{dispenseService})}]
START TRANSACTION;
SELECT status, patient_id FROM prescriptions WHERE prescription_id = ? FOR UPDATE;
SELECT medicine_id, stock_quantity FROM medicines WHERE medicine_id = ? FOR UPDATE;  -- each, ascending id
SELECT batch_id, quantity, expiry_date FROM medicine_batches
WHERE medicine_id = ? AND quantity > 0 AND expiry_date >= CURDATE()
ORDER BY expiry_date, batch_id FOR UPDATE;                                          -- first expiry first out
UPDATE prescription_items SET dispensed = TRUE, dispensed_at = NOW()
WHERE item_id IN (?) AND dispensed = FALSE;              -- trigger after_prescription_dispense lowers the total
UPDATE medicine_batches SET quantity = quantity - ? WHERE batch_id = ?;                -- each batch taken
INSERT INTO prescription_item_batches (item_id, batch_id, quantity) VALUES (?, ?, ?);
INSERT INTO pharmacy_stock_movements (medicine_id, movement_type, quantity, reason, reference_id, performed_by, batch_id)
VALUES (?, 'DISPENSE', ?, 'Dispense Prescription', ?, ?, ?);
UPDATE prescriptions SET status = 'DISPENSED' WHERE prescription_id = ?;
COMMIT;
\end{lstlisting}
The stress test runs 300 rounds of competing prescribing and dispensing on the same two medicines in opposite
orders, and 300 rounds of mixed writers: 0 deadlocks, and every medicine's total equals the sum of its batches.
Bed allocation locks the chosen bed and then the doctor (\tbl{FOR UPDATE}), so two clerks cannot give the same
last ICU bed to two patients.
